% !TeX root = ../papers/nc-p3.tex
\section{最小頂点次元四の均衡性と一隅五行列判定}
\label{sec:single-five-matrix-corner}
\begin{theorem}\label{thm:single-five-matrix-corner}
四辺有限な任意の標準正則窓で、有限次元非頂点単純のある頂点次元が4なら
全ての頂点次元は4である。これにはPI仮定を要しない。
従って四辺有限な標準有限Frobenius候補が正則かつ一つの指定隅で五行五列行列の恒等式を
全て満たすための必要十分条件も、定理\ref{thm:four-matrix-all-regular}の32次までの判定と
一般隅の指数検査である。
\end{theorem}
\begin{proof}[五次元の矢像での行数上界]
最小次元3以下は定理\ref{thm:generic-corner-small-simples}で済んでいる。
全位置が4以上として前節の $A_s:V\to E$ を用い、$m=\dim E$ と書く。
$s$ の空間は次元4以上で、各非零 $s$ で階数2以上、全像は $E$ を張る。
支持次元 $k=0,1,2,3,4$ の関係の次元上界を $\beta_k=0,0,1,3,6$ とする。
$A_s$ を選んで既存の支持の和を $t$ 次元増やせば、関係像の和は少なくとも
$6-\beta_{4-t}$ 増える。既存支持への逆像が余次元 $t$ だからである。
この増分下界は $t=0,1,2,3,4$ で $0,3,5,6,6$、逐次差は減少する。
従って $m=4q+r,0\le r<4$ なら関係像の和の次元は少なくとも $6q+6-\beta_{4-r}$ である。
$m=4$ では前節の証明で9以上である。入力次元が4より大きくても、
一般階数2の場合はどの四次元部分空間もそこで排除した定階数空間になってしまう。

$m=5$ では10以上であることを示す。9以下と仮定する。一般階数4ならそれらの像超平面は動く。
一般の異なる二超平面で関係像は各六次元で、その共通部分は三次元以上である。
支持の交わりの逆像は三平面だからその部分の関係はちょうど三次元である。
固定した一般の $A_s$ においてこのような逆像の候補は有限個である。
三次元の関係を支える超平面の方程式は辺の評価階数3の点だからである。
従って一般の像超平面は固定した三空間 $F\subset E$ を含む。
各一般 $s$ で $A_s^{-1}F$ も同じ有限性から一つの固定超平面 $K\subset V$ である。
$A_s(K)\subset F$ は全 $s$ に延び、$V/K\to E/F$ は二成分の一次式である。
非零な $s$ で両方零にでき、その $A_s$ の関係像は次元5以上で $F$ に支持される。
一方、元の一般二像の和は既に九次元で、$F$ に支持される部分は三次元しかない。
商の異なる二直線への成分が相殺できず、それぞれで商に消えるのは同じ三次元の共通部分だからである。

一般階数3なら、一般の像三空間の対は二次元で交わる必要がある。
交わりが一次元以下なら逆像は二次元以下なので二つの六次元関係像は高々一次元で交わる。
二次元交叉を持つ三空間族は、固定二空間 $F$ を含むか固定四空間に入る。
実際相異なる二元の共通部分を含まない第三元はその和四空間に入り、
そのような第三元があれば固定二空間を含む他の元も四空間の外へ出られない。
全像が張るから前者である。一般 $A_s$ で $F$ の逆像はやはり三次元支持の三関係を含む
固定超平面 $K$ で、$V/K\to E/F$ は三つの一次式である。
入力次元4以上より非零な共通零点を取れ、先と同じ五対三の矛盾となる。
一般階数2なら二像が交わらなければその関係像の和は既に10以上である。
全てが交わる像二平面の族は固定直線を含むか三空間に入り、後者は不可能。
前者では固定直線の商で独立な三方向を順に取る。最初に5以上、次の二回で少なくとも
3ずつ増すので9を超える。

$m=4$ の下界も10に上げられる。可逆元 $A$ があり関係像の和が9以下とする。
階数2の元があれば前節の下界は10なので、全特異元は階数3である。
その像超平面は $A$ で引き戻すと三関係を支え、従って候補は有限個である。
特異元の軌跡は行列式超曲面で、固定超平面に像が入る条件は入力の線型条件である。
ゆえに少なくとも射影平面の入力で像は一定となり、全点で階数3である。
しかしその平面で $4\mathcal O\to3\mathcal O(1)$ が全射なら核線束は
$c_2((1+H)^{-3})=6H^2$ を持つことになり矛盾する。
一般階数3なら、固定した一般像と他の像との交わりは二平面で、その逆像に三関係がなければ
和が10以上となる。五次元の一般階数4の場合と同じ有限性・鉛筆の議論で、
固定二平面 $F$、固定超平面 $K$、および $F$ に支持される少なくとも五次元の関係像を得て矛盾する。
一般階数2は前節で除外した。

以上から行解空間の上界 $L(m)$ を $L(4)=6,L(5)=10$、$m\ge6$ では
$\min(3(m-1),4m-6q-6+\beta_{4-r})$ と取れる。第一の上界は階数一行列の不在による。
行・列両方に適用する。
\end{proof}

\begin{proof}[二端の普遍拡張と複数の中間頂点]
非頂点単純 $S$ とその一般の異なる倍率 $T$ を取り、次元を $d_j$ とする。
$a_j=\dim\operatorname{Ext}^1(S,k_j)$、$b_j=\dim\operatorname{Ext}^1(k_j,S)$ と書く。
$h_j=2d_j-4(d_{j-1}+d_{j+1})+6d_{j+2}$、$H_j=h_j+a_j+b_j$、
$E_2=\sum d_jh_j$ とする。既出のEuler式と一般倍率分離により $H_j$ は頂点との第二Ext次元、
$E_2$ は二端間の第二Ext次元で、$E_2\ge\sum a_j b_j$ である。
前半より $a_j\le L(d_{j-1})-d_j$、$b_j\le L(d_{j+1})-d_j$。

相異なる $i,k$ で $A=a_i,B=b_k$ とし、いずれかを零としてもよい。
頭 $S$ の下に $k_i^A$ を置く普遍拡張を $U$、socle $T$ の上に $k_k^B$ を置くものを $W$ とする。
任意の $j$ について $\operatorname{Ext}^1(U,k_j)$ と $\operatorname{Ext}^1(k_j,W)$ の次元下界は
\[
 u_j=a_j+\mathbf1_{j=i+1}(4A-H_j)_+,\qquad
 v_j=b_j+\mathbf1_{j=k-1}(4B-H_j)_+
\]
と取れるが、$j=i$ で $A>0$ のときの第一項 $a_i$、$j=k$ で $B>0$ のときの
第二項 $b_k$ は、それぞれ普遍拡張の連結写像で消える。この二位置では該当項を零とする。
これはHomから始まる二つの長完全列で、他の頂点にはその減少がない。

$A,B$ の正の側の $i,k$ とは異なる頂点 $e$ を一つ選ぶと、
\[
 \sum_{j\ne e}u_jv_j
 \ \le\ E_2+A(h_i+a_i)+B(h_k+b_k)+6AB\mathbf1_{k=i+2}.
\]
左辺は $k_j$ を中間とするYoneda積の直和による。積の核の非零テンソルを階数分解して、
頭 $U$、中間がそれら頂点の和、下が $W$ の三層加群を作る。
$W$ で割った加群は $S$ のブロックで生成される。実際 $U$ が頭 $S$ で生成され、
その持上げが中間の和全体を生成しなければ、ある単純への押出しで第一拡張が分裂して
階数分解の独立性に反する。全加群で $S$ のブロックが生成する部分は $W$ に非零に交わる。
交わらなければ $W$ で割った加群への切断になり、中間に制限した下側の拡張も分裂してしまう。
$W$ のsocleは普遍性により $T$ だけなので、$S,T$ 間に非零矩形ができる。
他の全因子は $e$ を避けるから補題\ref{lem:generic-corner-separation}に反する。
従って上の直和は $\operatorname{Ext}^2(U,W)$ に単射となる。
右辺では、二つの完全列により $\operatorname{Ext}^2(S,T)$ の像を除く次元は
$AH_i+BH_k+6AB\mathbf1_{k=i+2}$ 以下である。その像では少なくとも
$Ab_i+a_kB$ 次元を失う。両者は普遍拡張で殺されるそれぞれ $k_i,k_k$ を通る積であり、
相異なる頂点を経る積の直和の単射性により独立だからである。
片側を使わない場合も同じ議論でその重複度を零にする。
\end{proof}

\begin{proof}[六次元を通過する積の改良と全列挙]
定理\ref{thm:opposite-four-matrix-corners}の $i,j=i+1,k=i+2$ に対する最後の不等式も用いる。
今 $d_{i+3}=6$ とする。中間六空間を通る行・列解空間の次元を
$p=d_i+a_i,l=d_k+b_k$ とすると、$l>10$ のとき全ての列解との積の関係評価を消す行解は
零または階数3以上である。実際、階数2の行 $F$ があれば、列解と掛けた像は
二文字の片側二次元支持を持つ $R^\perp$ に入る。その次元は
$8-6+\dim(R\cap (\text{相補な二次元支持}))\le3$ である。
従って四空間 $N_0=\ker F$ を値域に持つ列解の空間 $X_0$ は次元8以上となる。

元の表現の実際の各行を $N_0$ に制限する。その転置全体を $P_s:V\to N_0^*$ とすると、
その像は $N_0^*$ を張り、写像の空間も少なくとも二次元である。
前者は矢の共通零ベクトルの不在、後者は $N_0$ の各非零ベクトルで矢像が
少なくとも二次元であることによる。非零な $P_s$ は階数2以上である。
もし階数1なら、この行と $X_0$ の積は固定された片因子を持つ $R^\perp$ の可分解元で、
少なくとも二次元を張る。実際 $X_0$ の一つの出力共変元への射影の核は三行四列で
階数1を含まず高々6次元である。すると辺の点グラフに直線ができる。
ここで本来の行は全ての列解を零化するため、これらの積が $R^\perp$ に入る。

この $P_s$ による関係像の和は少なくとも9次元である。前節の四行の場合で、
一つ可逆なら独立な二写像の鉛筆でよく、一般階数3なら二つの異なる像があればよい。
一般階数2でも、二像が交わらなければ関係像は合計10以上、全てが交われば
共通直線を持つ族からその商で独立な三方向を取り $5+3+3$ 以上となる。
像が固定三空間に入る可能性は全体を張ることに反する。
一方 $X_0$ はこの和を零化するので次元は $16-9=7$ 以下であり矛盾する。

従って階数2以下の四行六列行列の余次元8によって
$r_a=(p-8)_+$ ($l>10$)、その他では0を、前掲の定理の有効次元下界に取れる。
逆向きに $r_b=(l-8)_+$ ($p>10$)、その他では0を取る。
補題\ref{lem:generic-corner-separation}を使う三段構成は頂点 $i+3$ を避けるので、
PI仮定なしに同定理の「二端だけで超過」の不等式を使える。すなわち $u_jv_j$ は先の右辺から
$\max(A r_a,B r_b)$ を引いたもの以下である。

以上は全て整数の有限検査になる。$d_0=4$、向きを反転して $4\le d_1\le d_3\le L(4)$、
$4\le d_2\le\min(L(d_1),L(d_3))$ とする。各 $a,b$ を非負の行列上界内で動かし、
$H_j\ge0$、$\sum a_jb_j\le E_2$ と上の二種の不等式を検査する。
\path{computations/verify_single_five_matrix_corner.py} は均衡型 $(4,4,4,4)$ だけを残す。

最後に一指定隅が五行列の恒等式を満たす場合、そこが次元5なら中心多項式で頂点に直交し均衡する。
4以下も今の定理と従前の結果で均衡する。定理\ref{thm:five-matrix-finite-windows}の残りの証明で
五行列恒等式を使った箇所は三つのブロックの和が5を超える隅での非零積の排除である。
全単純が均衡なのでいずれも同じ指定隅で行える。
また一隅から折り畳みのPI性が従い、素商の抽出とtrace繊維の次元の議論に他隅の恒等式は要らない。
よって中心有限な合成整域で指定隅の指数は2または4、全隅の一般指数は一致し同じ有限判定に帰着する。
\end{proof}
